\scanpage{57}
\mspage{28}
se grouperaient donc \deleted{dans de} par paquets (en regardant les centralisateurs des ss-groupes images de $\Gamma^\circ$ par ses sections,\deleted{)} et en mettant ensemble des scindages qui ont des $\unsure{\mathrm{Zentr}}(\Gamma^\circ)$ conjugués (même s'ils ne sont eux-mêmes conjugués). Donc on \unsure{retrouve} encore une description de $X(K)$ (donc aussi des $X(K')$ pour $\forall$ ext. finie \inserted{$K'$} de $K$) en termes de l'extension $E_L$ de $E_K$ par $\Pi_{L/K}$, en $\unsure{=|\overline m}$ temps qu'une façon de reconstituer les $U=X\setminus\bar I$\ldots Donc en fait c'est la structure $E_L\longrightarrow E_K$ qui est la plus riche a priori, et de loin plus commode pour le genre $0$ et $1$, où la considération des $U$ de type $(g,\nu)$, $(2g+\nu\geq3)$ casse le groupe ``continu'' d'automorphismes\ldots La forme ``modélique'' de la conjecture précédente revient à la forme ``birationnelle'', quand on y précise cette dernière en disant que \unsure{tt}
\margin{56}
